\chapter{Bayesian Inference}

In the chapter on maximum likelihood we learned to point at the single value of a
parameter that makes the data look least surprising, and to call that value our estimate.
It is a sharp, powerful idea, and it has one quietly unsatisfying feature: it hands you a
number and walks away. Ask it ``how sure are you?'' and it can only mumble about standard
errors computed from a curvature it pretends to trust. Ask it ``what did you believe
\emph{before} you saw the data?'' and it shrugs --- it has no notion of belief at all,
only of likelihood.

Bayesian inference is the other half of statistics, the half that takes the word
\emph{belief} seriously and treats it as something you can do arithmetic with. The idea
is breathtakingly simple, and you have already met its engine: Bayes' theorem, the rule
for running the arrow of cause and effect backwards. In the diagnosis chapter we used it
to update from a positive test to a disease. Here we do exactly the same thing, but the
``hidden cause'' is now a \emph{parameter} --- the bias of a coin, the rate of a Poisson
process, the mean of a normal --- and the ``evidence'' is a whole dataset. The output is
not a single number but an entire probability distribution over the parameter: a complete,
quantitative description of everything the data has taught us, uncertainty
included.

That distribution is called the \emph{posterior}, and learning to compute it, read it,
summarize it, and trust it is the business of this chapter. We will start with the
three-part incantation --- prior, likelihood, posterior --- and watch a real belief
sharpen as coin flips arrive. We will discover the small miracle of \emph{conjugate
priors}, where the posterior comes out in the same family as the prior and the update is a
one-line bookkeeping step. We will work three of these in full --- Beta--Binomial,
Gamma--Poisson, Normal--Normal --- and plot the posterior tightening with every data
point. We will learn to draw \emph{credible intervals} that mean exactly what a confidence
interval is so often mistaken to mean. We will contrast the single-point \emph{MAP}
summary with the full posterior, set the Bayesian and frequentist worldviews fairly side
by side, and finish in the engine room of modern machine learning, where these ideas run
billions of times a day.

\chapterline

\section{The Three Words: Prior, Likelihood, Posterior}
\label{ch31:three-words}

Let us fix a concrete question and never let go of it. A friend hands you a thumbtack and
asks: when I toss this and let it land, what is the probability $\theta$ that it lands
point-up? You do not know $\theta$. A coin would be $\theta = 1/2$, but a thumbtack is
lopsided, and really it could be anything. You will toss it a number of times, count the
point-ups, and try to learn $\theta$.

The frequentist instinct is to treat $\theta$ as a fixed but unknown constant and estimate
it. The Bayesian instinct is different and, once you see it, hard to unsee: treat your
\emph{uncertainty about} $\theta$ as itself a probability distribution. Before any tosses,
you have some vague sense of what $\theta$ might be --- maybe centered near $1/2$, maybe
flat across $[0,1]$. That sense is a distribution, the \emph{prior} $p(\theta)$. After the
tosses, your sense has sharpened into a new distribution, the \emph{posterior}
$p(\theta \given \text{data})$. The whole of Bayesian inference is the bridge between
those two.

\begin{intuitionbox}
A prior is not a guess at the answer. It is a description of \emph{your ignorance},
written in the language of probability. ``I have no idea, it could be anything in
$[0,1]$'' is the uniform prior. ``It is probably near a half but I would not bet the
house'' is a gentle hump around $0.5$. ``I have tossed thousands of these tacks before and
they land up about $60\%$ of the time'' is a tight spike near $0.6$. Each is a legitimate
starting belief, and --- this is the point --- each will be \emph{dragged toward the
truth by data}. The disagreement between two reasonable priors shrinks as evidence
accumulates. The prior is where you start; the data decides where you end up.
\end{intuitionbox}

The bridge is Bayes' theorem, written now for a continuous parameter. Let the data be $D$
and the parameter $\theta$. Then
\begin{equation}
  \underbrace{p(\theta \given D)}_{\text{posterior}}
  \;=\;
  \frac{\overbrace{p(D \given \theta)}^{\text{likelihood}}\;\overbrace{p(\theta)}^{\text{prior}}}
       {\underbrace{p(D)}_{\text{evidence}}}.
  \label{ch31:bayes-continuous}
\end{equation}
Three of these four objects you already know intimately. The \emph{prior} $p(\theta)$ is
your starting belief. The \emph{likelihood} $p(D \given \theta)$ is the same function the
maximum-likelihood chapter built its whole world on: for each candidate value of $\theta$,
how probable is the data we actually saw? The \emph{posterior} $p(\theta \given D)$ is the
updated belief we are after. The fourth, the \emph{evidence} $p(D) = \int p(D \given
\theta)\, p(\theta)\, \dd\theta$, is just the number that makes the posterior integrate to
one. It does not depend on $\theta$ at all. That single fact --- that the denominator is a
constant in $\theta$ --- is the most liberating thing in the subject.

\begin{focusbox}
Because the evidence is constant in $\theta$, almost all Bayesian reasoning runs on the
\emph{proportionality} form of Bayes' theorem:
\[
  p(\theta \given D)\;\propto\; p(D \given \theta)\, p(\theta),
  \qquad\text{in words:}\qquad
  \text{posterior} \;\propto\; \text{likelihood} \times \text{prior}.
\]
Compute the right-hand side as a function of $\theta$, recognize what distribution it is
(or normalize it at the end), and you are done. The phrase ``posterior is proportional to
likelihood times prior'' is the single most useful sentence in this chapter. Say it until
it is reflex.
\end{focusbox}

\begin{whybox}
Why is the posterior \emph{proportional} to likelihood times prior, and not equal to it?
Because the product $p(D\given\theta)p(\theta)$ is a perfectly good shape but the wrong
height: its total area is $p(D)$, not $1$. Dividing by $p(D)$ rescales the area to $1$
without touching the shape. So the prior$\,\times\,$likelihood product already \emph{is}
the posterior up to a vertical rescaling, and the rescaling is forced --- it is the unique
constant that makes a valid density. This is why we can so often ignore the denominator,
do the algebra on the shape, and read off the answer by pattern-matching to a known
family. The normalizer was never carrying information about $\theta$; it was only ever
carrying the bill for making probabilities sum to one.
\end{whybox}

\subsection{A First Update by Hand}

Let us toss the tack three times and get \emph{up, up, down}. Suppose, to keep the first
example arithmetic-free, that you have decided $\theta$ can only be one of three values:
$0.3$, $0.5$, or $0.7$, and before tossing you considered them equally plausible, so the
prior puts mass $1/3$ on each. The likelihood of two ups and one down for a given $\theta$
is $\theta^2(1-\theta)$. Multiply prior by likelihood:

\begin{center}
\begin{tabular}{cccc}
\toprule
$\theta$ & prior $p(\theta)$ & likelihood $\theta^2(1-\theta)$ & prior$\times$lik. \\
\midrule
$0.3$ & $1/3$ & $0.063$ & $0.0210$ \\
$0.5$ & $1/3$ & $0.125$ & $0.0417$ \\
$0.7$ & $1/3$ & $0.147$ & $0.0490$ \\
\midrule
\multicolumn{3}{r}{sum (the evidence):} & $0.1117$ \\
\bottomrule
\end{tabular}
\end{center}

Now normalize by dividing each product by the sum $0.1117$. The posterior is
$\P(\theta=0.3\given D) = 0.188$, $\P(\theta=0.5\given D)=0.373$, and
$\P(\theta=0.7\given D)=0.439$. We started believing all three equally and now lean toward
$0.7$, with $0.5$ a close second and $0.3$ trailing (Figure~\ref{ch31:fig-discrete}) ---
exactly the ordering of the
likelihoods, because here the flat prior contributed nothing to the \emph{shape}. That is
the whole machine. Everything else in this chapter is the same machine run on a continuum
of $\theta$ values instead of three, where the sums become integrals and the bookkeeping
becomes algebra.

\begin{figure}[H]
\centering
\begin{tikzpicture}
\begin{axis}[
  width=11cm, height=5.6cm,
  ybar, bar width=14pt,
  enlarge x limits=0.25,
  xlabel={candidate value of $\theta$},
  ylabel={probability mass},
  symbolic x coords={0.3,0.5,0.7},
  xtick=data,
  ymin=0, ymax=0.62,
  legend pos=north west, legend cell align=left,
  nodes near coords, nodes near coords style={font=\scriptsize},
  every node near coord/.append style={/pgf/number format/fixed,
    /pgf/number format/precision=3},
]
\addplot[draw=gray!70!black, fill=gray!25] coordinates {(0.3,0.333) (0.5,0.333) (0.7,0.333)};
\addlegendentry{prior (uniform)}
\addplot[draw=distblue, fill=distblue!35] coordinates {(0.3,0.188) (0.5,0.373) (0.7,0.439)};
\addlegendentry{posterior after $D=\{$up, up, down$\}$}
\end{axis}
\end{tikzpicture}
\caption{The first update made visible. The flat prior puts mass $1/3$ on each candidate;
after seeing two ups and one down the posterior tilts toward larger $\theta$, ranking the
candidates in exactly the order of their likelihoods $\theta^2(1-\theta)$.}
\label{ch31:fig-discrete}
\end{figure}

\begin{historybox}
The reverend Thomas Bayes (c.\ 1701--1761) never published the theorem that bears his
name; it was found in his papers after his death and read to the Royal Society in 1763 by
his friend Richard Price. Bayes' own example was, in spirit, exactly our thumbtack: a ball
rolled onto a table at an unknown position, with later balls reported only as landing left
or right of it, and the question of what could be inferred about the first ball's place.
It was Pierre-Simon Laplace, working independently a few years later, who turned the idea
into a general tool and used it on everything from the masses of planets to the ratio of
male to female births in Paris. For most of the next century the method was simply called
``inverse probability,'' and it was the dominant mode of statistical reasoning. The name
``Bayesian'' --- and the fierce twentieth-century rivalry with the frequentist school we
will meet later --- came much later still.
\end{historybox}

\section{The Beta--Binomial Conjugate Pair}
\label{ch31:beta-binomial}

Now let $\theta \in [0,1]$ range over the whole interval. We toss the tack $n$ times,
independently, and see $k$ point-ups. The likelihood is the binomial probability, viewed
as a function of $\theta$:
\[
  p(D \given \theta) \;=\; \binom{n}{k}\theta^{k}(1-\theta)^{\,n-k}
  \;\propto\; \theta^{k}(1-\theta)^{\,n-k},
\]
where we dropped $\binom{n}{k}$ because it does not involve $\theta$ and so disappears into
the proportionality. We need a prior on $[0,1]$. We could pick any density there, but one
family is magic for this likelihood: the \emph{Beta} distribution. Recall
$\theta \sim \Bet(\alpha,\beta)$ has density
\[
  p(\theta) \;=\; \frac{1}{B(\alpha,\beta)}\,\theta^{\alpha-1}(1-\theta)^{\beta-1},
  \qquad 0 \le \theta \le 1,
\]
where $B(\alpha,\beta)$ is the constant that makes it integrate to one. Look at the shape:
$\theta^{\alpha-1}(1-\theta)^{\beta-1}$. It is \emph{the same algebraic form as the
binomial likelihood}, a power of $\theta$ times a power of $1-\theta$. Multiply them and
the powers simply add:
\begin{align*}
  p(\theta \given D)
  &\;\propto\; \underbrace{\theta^{k}(1-\theta)^{n-k}}_{\text{likelihood}}
      \cdot \underbrace{\theta^{\alpha-1}(1-\theta)^{\beta-1}}_{\text{prior}} \\[2pt]
  &\;=\; \theta^{\,(\alpha+k)-1}\,(1-\theta)^{\,(\beta + n - k)-1}.
\end{align*}
The right-hand side is, up to its normalizing constant, exactly the density of a
$\Bet(\alpha + k,\; \beta + n - k)$. We did not have to compute a single integral. The
posterior is Beta because the prior was Beta and the likelihood was binomial; we say the
Beta family is \emph{conjugate} to the binomial likelihood.

\begin{formalbox}
\textbf{Beta--Binomial update.} If the prior is $\theta \sim \Bet(\alpha,\beta)$ and we
observe $k$ successes in $n$ independent Bernoulli trials, then the posterior is
\begin{equation}
  \theta \given D \;\sim\; \Bet(\alpha + k,\; \beta + n - k).
  \label{ch31:beta-update}
\end{equation}
In words: \emph{add your successes to $\alpha$ and your failures to $\beta$.} That is the
entire computation.
\end{formalbox}

\begin{happeningbox}
The update rule \eqref{ch31:beta-update} tells you exactly what a Beta prior \emph{is}: a
record of imaginary past data. A $\Bet(\alpha,\beta)$ prior behaves as though, before any
real tosses, you had already seen $\alpha-1$ successes and $\beta-1$ failures. (The
$\Bet(1,1)$ prior, which is flat, corresponds to zero imaginary data: total openness.)
Seen this way, $\alpha$ and $\beta$ are \emph{pseudocounts}, and Bayesian updating is
nothing but continuing to count, starting from a head start you chose. Real successes and
imaginary successes are added in exactly the same way. This is why a strong prior (large
$\alpha,\beta$) resists the data --- it is outvoted only slowly --- while a weak prior is
overwhelmed almost at once. The prior's strength is literally measured in units of data.
\end{happeningbox}

\subsection{Watching the Posterior Tighten}

Suppose we begin with the flat $\Bet(1,1)$ prior --- complete openness about the tack ---
and the true bias is $\theta = 0.7$. We toss it and update after $5$, $20$, and $80$
tosses, getting close to $70\%$ ups each time (as close as the discreteness of small
samples allows). The posteriors are
$\Bet(1+4,\,1+1)=\Bet(5,2)$ after $5$ tosses ($4$ up, $1$ down),
$\Bet(1+14,\,1+6)=\Bet(15,7)$ after $20$ ($14$ up, $6$ down), and
$\Bet(1+56,\,1+24)=\Bet(57,25)$ after $80$ ($56$ up, $24$ down). Figure~\ref{ch31:fig-tighten} plots all three
against the flat prior. Each new batch of data pulls the hump toward $0.7$ and squeezes it
narrower: the belief is becoming both more accurate and more confident at the same time.

\begin{figure}[H]
\centering
\begin{tikzpicture}
\begin{axis}[
  width=12.5cm, height=7cm,
  xlabel={$\theta$ (probability of point-up)},
  ylabel={posterior density},
  domain=0.001:0.999, samples=160,
  xmin=0, xmax=1, ymin=0, ymax=9,
  legend pos=north west, legend cell align=left,
  every axis plot/.append style={very thick},
  axis lines=left,
]
% Each Beta density is written as exp(logC + (a-1)ln x + (b-1)ln(1-x)),
% with logC = ln Gamma(a+b) - ln Gamma(a) - ln Gamma(b) precomputed exactly.
% This log form is numerically robust (no overflow in the normalizer).
% prior Beta(1,1): density = 1
\addplot[gray, dashed] {1};
\addlegendentry{prior $\Bet(1,1)$}
% Beta(5,2): logC = 3.40120
\addplot[distblue] {exp(3.40120 + 4*ln(x) + 1*ln(1-x))};
\addlegendentry{after 5 tosses, $\Bet(5,2)$}
% Beta(15,7): logC = 13.60967
\addplot[quotientorange] {exp(13.60967 + 14*ln(x) + 6*ln(1-x))};
\addlegendentry{after 20 tosses, $\Bet(15,7)$}
% Beta(57,25): logC = 50.93005
\addplot[densityred] {exp(50.93005 + 56*ln(x) + 24*ln(1-x))};
\addlegendentry{after 80 tosses, $\Bet(57,25)$}
\draw[black, thin, dotted] (axis cs:0.7,0) -- (axis cs:0.7,9);
\node[anchor=south] at (axis cs:0.7,8.2) {\scriptsize truth $0.7$};
\end{axis}
\end{tikzpicture}
\caption{The posterior tightening as data arrives. Starting from a flat prior, each batch
of tosses (all roughly $70\%$ point-up) pulls the density toward the true value $0.7$ and
narrows it. More data means a taller, thinner posterior --- more confident and more
accurate at once. (Curves are drawn from closed-form Beta densities; exact heights are
illustrative.)}
\label{ch31:fig-tighten}
\end{figure}

The posterior mean of a $\Bet(a,b)$ is $a/(a+b)$. After $80$ tosses that is
$57/82 \approx 0.695$, gratifyingly close to the truth, and after only $5$ tosses it was
$5/7 \approx 0.714$ --- already in the right neighborhood but, as the wide curve warns, not
yet to be trusted. This is the great virtue of carrying the whole distribution rather than
a point: the \emph{width} of the posterior is a frank measure of how much you do not
yet know.

\begin{walkbox}
Notice a subtle thing about the posterior mean. We can write it as
\[
  \frac{\alpha + k}{\alpha + \beta + n}
  \;=\; \underbrace{\frac{\alpha+\beta}{\alpha+\beta+n}}_{w}\cdot
        \underbrace{\frac{\alpha}{\alpha+\beta}}_{\text{prior mean}}
   \;+\; \underbrace{\frac{n}{\alpha+\beta+n}}_{1-w}\cdot
        \underbrace{\frac{k}{n}}_{\text{data mean}}.
\]
The posterior mean is a \emph{weighted average} of the prior mean and the sample
proportion $k/n$, with the weight on the data being $n/(\alpha+\beta+n)$. When $n$ is small
the prior dominates; as $n\to\infty$ the weight on the data tends to $1$ and the Bayesian
estimate converges to the frequentist one. The two schools agree in the limit of abundant
data; they differ most exactly when data is scarce, which is precisely when being candid about
uncertainty matters most.
\end{walkbox}

\begin{example}[A/B testing a button color]
\label{ch31:ex-ab}
You change a checkout button from gray to green and show the green version to $40$ visitors;
$12$ buy. With no strong prior belief, take $\theta\sim\Bet(1,1)$ for the green
conversion rate. The posterior is $\Bet(1+12,\,1+28) = \Bet(13,29)$. Its mean is
$13/42 \approx 0.310$, and we can ask any question we like of the full distribution: the
probability that the true rate exceeds $0.25$ is $\int_{0.25}^{1} p(\theta\given D)\,
\dd\theta \approx 0.79$. If the old gray button is known to convert at $0.25$, we can say
``there is about a $79\%$ chance the green button is genuinely better,'' a sentence a
frequentist $p$-value cannot legally utter. That is what has made Bayesian A/B testing
popular in industry: the output is a direct probability about the thing you care about.
\end{example}

\section{The Gamma--Poisson Conjugate Pair}
\label{ch31:gamma-poisson}

Counts in fixed windows --- emails per hour, goals per match, mutations per genome --- are
the natural home of the Poisson distribution with rate $\lambda$. Suppose we observe counts
$x_1,\dots,x_n$ from $\Pois(\lambda)$ and want to learn $\lambda > 0$. The likelihood is
\[
  p(D \given \lambda) \;=\; \prod_{i=1}^{n}\frac{\lambda^{x_i}e^{-\lambda}}{x_i!}
  \;\propto\; \lambda^{\sum_i x_i}\, e^{-n\lambda},
\]
again dropping factors free of $\lambda$. We need a prior on the positive reals whose shape
matches $\lambda^{(\cdot)}e^{-(\cdot)\lambda}$. That is exactly the \emph{Gamma}
distribution: $\lambda\sim\Gam(\alpha,\beta)$ (shape $\alpha$, rate $\beta$) has density
\[
  p(\lambda) \;\propto\; \lambda^{\alpha-1}e^{-\beta\lambda},\qquad \lambda>0.
\]
Multiply prior by likelihood and, as before, the powers and the exponents simply combine:
\[
  p(\lambda\given D)\;\propto\;
  \lambda^{\sum_i x_i}e^{-n\lambda}\cdot\lambda^{\alpha-1}e^{-\beta\lambda}
  \;=\;\lambda^{\,(\alpha + \sum_i x_i)-1}\,e^{-(\beta + n)\lambda},
\]
which is the kernel of a $\Gam(\alpha + \sum_i x_i,\; \beta + n)$.

\begin{formalbox}
\textbf{Gamma--Poisson update.} With prior $\lambda\sim\Gam(\alpha,\beta)$ and observed
counts summing to $S=\sum_i x_i$ over $n$ windows,
\begin{equation}
  \lambda\given D \;\sim\; \Gam(\alpha + S,\; \beta + n).
  \label{ch31:gamma-update}
\end{equation}
Add the total count to the shape; add the number of observations to the rate.
\end{formalbox}

Here too the parameters read as pseudo-data: $\beta$ is an imaginary number of prior
observation-windows and $\alpha$ an imaginary number of prior events. The posterior mean is
$(\alpha + S)/(\beta + n)$, again a weighted blend of the prior mean $\alpha/\beta$ and the
data mean $S/n$.

\begin{example}[Goals per match]
\label{ch31:ex-goals}
A football analyst models a team's goals per match as $\Pois(\lambda)$. From long
experience she takes a $\Gam(2,1)$ prior --- prior mean $2$ goals, fairly broad. The team
then scores $3, 1, 2, 4, 0$ in five matches, so $S=10$, $n=5$. The posterior is
$\Gam(2+10,\,1+5)=\Gam(12,6)$, with mean $12/6 = 2.0$ goals per match and standard
deviation $\sqrt{12}/6 \approx 0.58$. Five matches later the team scores $2,2,3,1,2$,
adding $S'=10$, $n'=5$; updating again gives $\Gam(22,11)$, mean $2.0$, standard deviation
$\sqrt{22}/11\approx 0.43$. Same central estimate, tighter belief --- the posterior has
narrowed because evidence has doubled (Figure~\ref{ch31:fig-gammapoisson}). Notice we updated in two steps and got the same
answer as one combined step would: \emph{Bayesian updating is associative}, a property we
examine next.
\end{example}

\begin{figure}[H]
\centering
\begin{tikzpicture}
\begin{axis}[
  width=12.5cm, height=6.4cm,
  xlabel={$\lambda$ (goals per match)},
  ylabel={density},
  domain=0.001:5, samples=160,
  xmin=0, xmax=5, ymin=0, ymax=1.1,
  legend pos=north east, legend cell align=left,
  every axis plot/.append style={very thick},
  axis lines=left,
]
% Gamma(a, rate b) density = exp(logC + (a-1) ln x - b x), logC = a ln b - lgamma(a).
% prior Gamma(2,1): logC = 0
\addplot[gray, dashed] {exp(0 + 1*ln(x) - 1*x)};
\addlegendentry{prior $\Gam(2,1)$}
% posterior after 5 matches Gamma(12,6): logC = 3.99881
\addplot[quotientorange] {exp(3.99881 + 11*ln(x) - 6*x)};
\addlegendentry{after 5 matches, $\Gam(12,6)$}
% posterior after 10 matches Gamma(22,11): logC = 7.37356
\addplot[densityred] {exp(7.37356 + 21*ln(x) - 11*x)};
\addlegendentry{after 10 matches, $\Gam(22,11)$}
\draw[black, thin, dotted] (axis cs:2,0) -- (axis cs:2,1.0);
\node[anchor=south] at (axis cs:2,0.97) {\scriptsize mean $2.0$};
\end{axis}
\end{tikzpicture}
\caption{Gamma--Poisson updating for the goals example. The broad $\Gam(2,1)$ prior
sharpens to $\Gam(12,6)$ after five matches and to $\Gam(22,11)$ after ten. The mean stays
at $2.0$ goals per match while the density grows taller and narrower: doubling the evidence
roughly halves the variance, so the standard deviation falls from $0.58$ to $0.43$.}
\label{ch31:fig-gammapoisson}
\end{figure}

\begin{aestheticbox}
The fact that you can feed data in one observation at a time, or all at once, or in any
grouping you like, and always reach the same posterior, is not a lucky coincidence. It is
forced by the structure of \eqref{ch31:bayes-continuous}: yesterday's posterior is exactly
today's prior. Belief flows forward through time like a ledger that only ever gets
appended to. This is why Bayesian methods are the natural language of \emph{online
learning} and \emph{streaming data} --- a spam filter, a recommendation engine, a robot's
map of its surroundings. None of them can afford to re-process all of history at every
step; all of them can afford to update a posterior. The conjugate families make that update
a handful of additions.
\end{aestheticbox}

\section{The Normal--Normal Conjugate Pair}
\label{ch31:normal-normal}

The third classic conjugate pair governs measurement. Suppose we measure a quantity whose
true value is $\mu$, with each measurement $x_i \sim \Normal(\mu, \sigma^2)$ where the
measurement noise $\sigma^2$ is known and only the mean $\mu$ is unknown. Take a normal
prior on the mean, $\mu \sim \Normal(\mu_0, \tau_0^2)$, encoding a prior best guess $\mu_0$
with prior uncertainty $\tau_0^2$.

The algebra is slightly heavier because of the exponentials, but the punchline is clean and
worth deriving. With $n$ observations of sample mean $\bar x$, the likelihood is (dropping
constants)
\[
  p(D\given\mu)\;\propto\;\exp\!\left(-\frac{n(\mu-\bar x)^2}{2\sigma^2}\right),
\]
and the prior is $\propto \exp\!\big(-(\mu-\mu_0)^2/(2\tau_0^2)\big)$. Their product is the
exponential of a sum of two quadratics in $\mu$, which is again a quadratic in $\mu$, which
is again (the kernel of) a normal. Completing the square gives the posterior parameters.

\begin{formalbox}
\textbf{Normal--Normal update (known variance).} With prior $\mu\sim\Normal(\mu_0,\tau_0^2)$
and $n$ observations of mean $\bar x$ from $\Normal(\mu,\sigma^2)$, the posterior is
$\mu\given D \sim \Normal(\mu_n, \tau_n^2)$ where it is cleanest in terms of
\emph{precision} (the reciprocal of variance):
\begin{equation}
  \frac{1}{\tau_n^2} = \frac{1}{\tau_0^2} + \frac{n}{\sigma^2},
  \qquad
  \mu_n = \tau_n^2\!\left(\frac{\mu_0}{\tau_0^2} + \frac{n\,\bar x}{\sigma^2}\right).
  \label{ch31:normal-update}
\end{equation}
\end{formalbox}

Stare at \eqref{ch31:normal-update} until it speaks. \emph{Precisions add.} Your prior
carries precision $1/\tau_0^2$; each of the $n$ data points carries precision
$1/\sigma^2$; the posterior precision is simply their sum. And the posterior mean $\mu_n$
is a precision-weighted average of the prior mean and the data mean --- each source pulls
the answer toward itself in proportion to how precise it is. A confident prior (large
$1/\tau_0^2$) and a noisy measurement (small $1/\sigma^2$) leave the posterior near the
prior; abundant precise data overwhelm the prior. It is the same weighted-average story as
Beta--Binomial and Gamma--Poisson, now told in the currency of precision.

\begin{whybox}
Why do \emph{precisions} add rather than variances? Because in the exponent of a Gaussian,
the coefficient of $\mu^2$ is $-\tfrac12 \times (\text{precision})$, and multiplying two
Gaussians adds the exponents, hence adds those coefficients, hence adds precisions. Variance
is the inverse, so it does the opposite of adding --- combining two equally precise sources
\emph{halves} the variance. This is the deep reason the standard error of a mean falls like
$\sigma/\sqrt n$: each independent observation contributes one unit of precision, $n$ of
them give precision $n/\sigma^2$, and the posterior standard deviation is the reciprocal
square root. The familiar $1/\sqrt n$ is precision-addition wearing a disguise.
\end{whybox}

\begin{figure}[H]
\centering
\begin{tikzpicture}
\begin{axis}[
  width=12.5cm, height=6.4cm,
  xlabel={$\mu$ (true value being measured)},
  ylabel={density},
  domain=-2:6, samples=180,
  xmin=-2, xmax=6, ymin=0, ymax=1.5,
  legend pos=north east, legend cell align=left,
  every axis plot/.append style={very thick},
  axis lines=left,
]
% prior N(0, 4): sd=2
\addplot[gray, dashed] {1/(2*sqrt(2*pi))*exp(-(x-0)^2/(2*4))};
\addlegendentry{prior $\Normal(0,4)$}
% likelihood centered at xbar=3, n=4, sigma^2=1 => effective precision 4, var 0.25, sd 0.5
\addplot[quotientorange] {1/(0.5*sqrt(2*pi))*exp(-(x-3)^2/(2*0.25))};
\addlegendentry{data ($\bar x=3$, $n=4$)}
% posterior: prec = 1/4 + 4 = 4.25, var=0.2353, sd=0.485; mean = 0.2353*(0/4 + 12) = 2.823
\addplot[densityred] {1/(0.4851*sqrt(2*pi))*exp(-(x-2.823)^2/(2*0.2353))};
\addlegendentry{posterior $\Normal(2.82,0.235)$}
\end{axis}
\end{tikzpicture}
\caption{Normal--Normal updating. A broad prior centered at $0$ meets data centered at
$\bar x = 3$. The posterior (red) sits between them --- pulled most of the way toward the
data because the data are far more precise --- and is narrower than either input, since
precisions add. The compromise is automatic and quantitative.}
\label{ch31:fig-normalnormal}
\end{figure}

\begin{connectionbox}
The precision-weighted average in \eqref{ch31:normal-update} is the one-dimensional seed of
an idea that grows into the \emph{Kalman filter}, the algorithm steering spacecraft,
phones, and self-driving cars. There the ``prior'' is your prediction of where a moving
object is, the ``data'' is a noisy sensor reading, and the posterior is the fused estimate
--- exactly a precision-weighted blend, run once per time step, in many dimensions with the
scalar variances replaced by covariance matrices. If you want to see where the linear
algebra enters, the companion linear-algebra reader's treatment of positive-definite
matrices and quadratic forms is the language in which the multivariate version is written:
the exponent $-(\mu-\mu_0)^2/(2\tau_0^2)$ becomes $-\tfrac12(\boldsymbol\mu -
\boldsymbol\mu_0)^\top \Sigma^{-1}(\boldsymbol\mu - \boldsymbol\mu_0)$, and ``precisions
add'' becomes ``inverse-covariance matrices add.''
\end{connectionbox}

\section{Summarizing a Posterior: Credible Intervals and MAP}
\label{ch31:summaries}

The posterior is the complete answer, but humans want a headline. Two summaries dominate.

\subsection{Credible Intervals}

A \emph{credible interval} is an interval that contains the parameter with a stated
posterior probability. A $95\%$ credible interval $[L,U]$ is any pair with
\[
  \int_{L}^{U} p(\theta\given D)\,\dd\theta = 0.95.
\]
Read that integral literally: \emph{given the data, there is a $95\%$ probability that
$\theta$ lies in $[L,U]$.} That is the sentence everyone wishes a confidence interval
meant. It is the sentence a credible interval actually means.

\begin{trapbox}
A frequentist $95\%$ \emph{confidence} interval does \emph{not} say there is a $95\%$
probability the parameter lies inside it. To a frequentist the parameter is a fixed
constant, so it is either in the interval or not --- no probability about it. The $95\%$
refers to the \emph{procedure}: if you repeated the whole experiment many times, $95\%$ of
the intervals so constructed would cover the true value. The Bayesian credible interval
makes the natural-language statement directly, at the cost of having required a prior. Do
not mix the two interpretations; the difference is the entire frequentist--Bayesian
fault line in miniature.
\end{trapbox}

There are two common ways to pin down the interval, because the equation above has many
solutions. The \emph{equal-tailed} interval chops $2.5\%$ off each end, taking $[L,U]$ to be
the $2.5$th and $97.5$th posterior percentiles --- easy and the usual default. The
\emph{highest posterior density} (HPD) interval instead takes the shortest interval with
$95\%$ mass, equivalently the set of $\theta$ whose posterior density exceeds some
threshold. For a symmetric unimodal posterior the two coincide; for a skewed one the HPD is
shorter and always contains the mode, which the equal-tailed one need not.

\begin{example}[Credible interval for the A/B test]
\label{ch31:ex-credible}
Return to the green button of Example~\ref{ch31:ex-ab}, posterior $\Bet(13,29)$. Its
equal-tailed $95\%$ credible interval runs from about $0.181$ to $0.455$ (the $2.5$th and
$97.5$th percentiles of that Beta, read from software or tables). So we report: ``given the
data, the green conversion rate is $95\%$ likely to lie between roughly $18\%$ and $46\%$,
with posterior mean $31\%$.'' Compact, accurate, and directly interpretable.
\end{example}

\subsection{MAP versus the Full Posterior}

The \emph{maximum a posteriori} estimate is the single most probable value of $\theta$
under the posterior --- the location of the peak:
\[
  \hat\theta_{\text{MAP}} \;=\; \arg\max_\theta\; p(\theta\given D)
  \;=\; \arg\max_\theta\; p(D\given\theta)\,p(\theta).
\]
The second equality (we dropped the constant evidence) reveals MAP as maximum likelihood
\emph{with a prior attached}: maximize likelihood times prior instead of likelihood alone.
Take a flat prior and MAP collapses exactly to the MLE. So MAP is the Bayesian's bridge to
the familiar point estimate, and the prior term acts as a \emph{regularizer} that pulls the
estimate away from extreme values the data alone might suggest.

\begin{happeningbox}
MAP throws away almost everything. It reports the tallest point of the posterior and
discards the width, the skew, the tails --- all the information about how \emph{sure} you
are. For a sharply peaked posterior that is a fine summary; for a broad or skewed or
multimodal one it can be badly misleading, even reporting a value of essentially zero
density elsewhere as if it were the answer. A notorious case: if a posterior has two
separated peaks, MAP names one and pretends the other does not exist. The full posterior is
the object that actually contains your knowledge; MAP, the posterior mean, and credible
intervals are all just lossy summaries of it. Whenever you can carry the whole thing, do.
\end{happeningbox}

\begin{walkbox}
For the green-button posterior $\Bet(13,29)$, compare the summaries. The MAP (mode of a
Beta is $(\alpha-1)/(\alpha+\beta-2)$) is $12/40 = 0.300$. The posterior mean is
$13/42 \approx 0.310$. They differ because the Beta is mildly right-skewed --- the mode
sits a touch left of the mean, with a longer tail to the right. Which should you report?
If you must give one number and downstream calculations average over many such estimates,
the posterior mean is usually the better choice (it minimizes expected squared error); if
you want the single most plausible value, MAP. But the most informative report is the interval from
Example~\ref{ch31:ex-credible}, which contains both and shows the spread.
\end{walkbox}

\section{Bayesian versus Frequentist: An Even-Handed Comparison}
\label{ch31:vs}

It would be a disservice to present Bayesian inference as simply correct and frequentist
inference as confused. They answer different questions, and the practitioner's job is to
know which question is being asked.

\begin{center}
\renewcommand{\arraystretch}{1.3}
\begin{tabular}{p{4.3cm} p{5.0cm} p{5.0cm}}
\toprule
 & \textbf{Frequentist} & \textbf{Bayesian} \\
\midrule
What is $\theta$? & A fixed unknown constant. & A random variable with a distribution. \\
What is random? & The data (over hypothetical repetitions). & Your belief about $\theta$, given the one dataset. \\
Need a prior? & No. & Yes --- this is the cost and the power. \\
Main output & Point estimate $+$ confidence interval; $p$-values. & Full posterior distribution. \\
``$95\%$ interval'' means & $95\%$ of such intervals cover $\theta$. & $\P(\theta \in \text{interval}\given\text{data})=0.95$. \\
Strength & No prior needed; strong long-run guarantees. & Direct probability statements; uses prior knowledge; clean with small data. \\
\bottomrule
\end{tabular}
\end{center}

\begin{paradoxbox}
\textbf{The same data, two different intervals --- and the trouble with optional stopping.}
Here is a genuinely unsettling fact. A frequentist analysis can depend on your
\emph{intentions}, not just your data. Suppose you flip a coin and see $9$ heads in $12$
flips. A frequentist who planned to flip exactly $12$ times computes one $p$-value; a
frequentist who planned to flip ``until I get $3$ tails'' (and happened to stop at the
twelfth flip) computes a \emph{different} $p$-value from the \emph{same observed sequence},
because the sample space of hypothetical repetitions is different. The data are identical;
the conclusion shifts with the experimenter's private stopping rule. The Bayesian is
untroubled: the likelihood of the observed flips is $\theta^9(1-\theta)^3$ either way, the
stopping rule contributes only a constant that the proportionality washes out, and the
posterior is identical. This is the \emph{likelihood principle}, and Bayesian inference
obeys it automatically. To some this is the decisive argument for the Bayesian view; to
others it is a curiosity. Either way it shows the two frameworks are genuinely, not
cosmetically, different.
\end{paradoxbox}

\begin{historybox}
The twentieth century saw a long cold war between the camps. Ronald Fisher, Jerzy Neyman,
and Egon Pearson built the frequentist edifice --- $p$-values, confidence intervals,
hypothesis tests --- partly in explicit reaction against ``inverse probability,'' which
they distrusted for its reliance on subjective priors. Harold Jeffreys kept the Bayesian
flame through mid-century; Leonard Savage and Bruno de Finetti gave it rigorous decision-
and belief-theoretic foundations. For decades Bayesian methods were admired in theory and
avoided in practice, for one brutally simple reason: the evidence integral $\int p(D\given
\theta)p(\theta)\,\dd\theta$ is usually impossible to do by hand outside the conjugate
families. What broke the stalemate was not a new theorem but cheap computers and the
\emph{Markov chain Monte Carlo} revolution of the late 1980s and early 1990s, which let us
\emph{sample} from any posterior without ever computing that integral. Today the
distinction in practice is less a war than a toolkit; most working statisticians use both.
\end{historybox}

\begin{aestheticbox}
The conjugate families of this chapter are beautiful precisely because they sidestep the
hard integral --- the posterior pops out in closed form. But they are a small, special
island. The reason Bayesian inference is now everywhere is that we learned to leave the
island: when the posterior has no closed form, draw samples from it with Markov chain Monte
Carlo (the companion to this chapter on Monte Carlo methods builds exactly this machinery)
and summarize the samples. Conjugacy is the place to build intuition; sampling is the place
where modern Bayesian work actually happens.
\end{aestheticbox}

\section{Bayesian Ideas in Machine Learning}
\label{ch31:ml}

The vocabulary of this chapter is the hidden skeleton of a great deal of machine learning.
Three of these connections deserve to be spelled out.

\begin{applicationbox}
\textbf{Regularization is a prior in disguise.} Training a model by minimizing a loss plus
a penalty $\lambda\norm{w}^2$ (``ridge'' or weight decay) is \emph{exactly} MAP estimation
with a Gaussian prior $w\sim\Normal(0,\,\sigma^2 I)$ on the weights. Take the log of
posterior $\propto$ likelihood $\times$ prior: the log-likelihood is the usual fit term,
and the log of a Gaussian prior is $-\norm{w}^2/(2\sigma^2)$ plus a constant --- precisely
the squared penalty, with $\lambda = 1/(2\sigma^2)$. Likewise an $L_1$ penalty (``lasso'')
is MAP under a Laplace prior, whose sharp peak at zero is what drives weights exactly to
zero. So every time a practitioner ``adds regularization to prevent overfitting,'' they
are, whether they know it or not, declaring a prior belief that the weights are small. The
strength $\lambda$ is the prior's confidence. Overfitting, in this light, is what happens
when you let a flexible model believe the data with no prior skepticism at all.
\end{applicationbox}

\begin{applicationbox}
\textbf{Na\"ive Bayes and posterior class probabilities.} A spam filter, or any
probabilistic classifier, computes $\P(\text{class}\given\text{features}) \propto
\P(\text{features}\given\text{class})\,\P(\text{class})$ --- prior times likelihood,
exactly \eqref{ch31:bayes-continuous} with the class playing the role of $\theta$. The
class prior $\P(\text{spam})$ is your base rate; the likelihood is the word model; the
posterior is the calibrated probability that the message is spam. The same shape powers
medical-diagnosis nets and topic models. When a modern neural network ends in a softmax and
you read its outputs as ``class probabilities,'' you are reading an approximate posterior.
\end{applicationbox}

\begin{applicationbox}
\textbf{Calibrated uncertainty in deep learning.} A standard neural network gives a single
weight vector --- a point estimate, the MAP of an implicit posterior. \emph{Bayesian deep
learning} instead tries to carry a distribution over weights, so the network can say ``I do
not know'' on inputs unlike anything it trained on, by reporting a wide posterior
predictive. This matters enormously in self-driving cars and medical AI, where a confident
wrong answer is far more dangerous than a candid ``uncertain.'' Exact posteriors over
millions of weights are hopeless, so the field approximates --- variational inference,
Monte Carlo dropout, deep ensembles --- but the goal is pure Bayes: replace the single
best guess with a calibrated distribution of guesses.
\end{applicationbox}

\begin{connectionbox}
There is also a clean bridge to information theory, the subject of a later chapter. The
amount a single data point changes your belief --- how far the posterior moves from the
prior --- is naturally measured by the Kullback--Leibler divergence $\KL(p(\theta\given
D)\,\|\,p(\theta))$, the ``information gain.'' Experimental design can then be framed as
choosing the experiment expected to yield the largest information gain, a principle called
\emph{Bayesian optimal design}. Belief updating and information are two views of one thing:
to learn is to reduce uncertainty, and uncertainty is measured in bits.
\end{connectionbox}

\section{The Posterior Predictive: Forecasting the Next Observation}
\label{ch31:predictive}

We have been learning about $\theta$, but often what we truly want is to predict the
\emph{next data point} $x_{\text{new}}$. The Bayesian answer does not plug in a single
estimate of $\theta$; it averages the prediction over the entire posterior, weighting each
value of $\theta$ by how plausible the data have made it. This is the \emph{posterior
predictive distribution}:
\begin{equation}
  p(x_{\text{new}}\given D) \;=\; \int p(x_{\text{new}}\given\theta)\,
       p(\theta\given D)\,\dd\theta.
  \label{ch31:postpred}
\end{equation}
Read it as a vote: every candidate $\theta$ predicts the future, and we tally the
predictions weighted by posterior belief. This automatically folds parameter uncertainty
into the forecast --- the predictive is wider than if we had pretended to know $\theta$
exactly, which is exactly right, because we do not.

\begin{example}[Will the next tack land up?]
\label{ch31:ex-postpred}
After the $80$ tosses with posterior $\Bet(57,25)$, what is the probability the next toss
lands point-up? Plug the Bernoulli likelihood $p(x_{\text{new}}=1\given\theta)=\theta$ into
\eqref{ch31:postpred}:
\[
  \P(x_{\text{new}}=1\given D) = \int_0^1 \theta\,p(\theta\given D)\,\dd\theta
  = \E[\theta\given D] = \frac{57}{57+25} = \frac{57}{82}\approx 0.695.
\]
The predictive probability is just the posterior \emph{mean} of $\theta$ --- a tidy result
special to the Bernoulli case. Note we did \emph{not} use the MAP value $56/80 = 0.700$;
averaging over the posterior, not maximizing, is the principled forecast. With few data the
gap between the two matters; with much data it vanishes.
\end{example}

\begin{labbox}
Try this with pencil and a tiny prior. Start with $\Bet(2,2)$ (a gentle belief that a coin
is fair) and feed it the sequence H, H, H, T one flip at a time, writing the posterior
after each flip using the rule ``add a success to $\alpha$, a failure to $\beta$.'' You
should pass through $\Bet(3,2),\Bet(4,2),\Bet(5,2),\Bet(5,3)$. After each step compute the
posterior-predictive probability of the \emph{next} flip being heads, $\alpha/(\alpha+\beta)$,
and watch it climb $0.6, 0.667, 0.714, 0.625$ as the evidence accumulates and then the lone
tail pulls it back. Feel how the gentle prior keeps the early estimates from lurching all
the way to $1$ after three straight heads --- that restraint is the prior doing its job.
\end{labbox}

\section{Conclusion}
\label{ch31:conclusion}

Bayesian inference is one idea applied relentlessly: \emph{posterior $\propto$ likelihood
$\times$ prior}. The prior states what you believed before; the likelihood states what the
data say; the posterior is the disciplined combination, a full distribution over the
unknown rather than a lonely point. We saw that for the three great conjugate pairs ---
Beta--Binomial, Gamma--Poisson, Normal--Normal --- the combination is a one-line update of
counts or precisions, and we watched the posterior tighten toward the truth as data
arrived, its narrowing width a true accounting of remaining doubt. We learned to read
off credible intervals that mean what people wish confidence intervals meant, to summarize
with MAP while remembering how much MAP discards, and to predict the future by averaging
over the posterior rather than betting on its peak.

The deepest lesson is one of attitude. A frequentist asks, ``what procedure has good
long-run behavior?'' A Bayesian asks, ``what should I believe, right now, given exactly
what I have seen?'' Both are legitimate; the second happens to be the question we ask in
ordinary life, and it is the one that, scaled up on fast computers, now quietly powers
spam filters, recommendation engines, robot navigation, and the calibrated uncertainty of
modern AI. You began this book learning to compute the probability of outcomes. You end
this chapter doing the inverse and far more human thing: computing, and rationally
updating, your beliefs about the hidden causes behind what you see.

\section{Exercises}
\label{ch31:exercises}

\begin{exercise}[Discrete update]
A box contains a coin that is either fair ($\theta = 0.5$) or two-headed ($\theta=1$), with
prior probabilities $0.9$ and $0.1$. You flip it once and get heads. Find the posterior
probability that the coin is two-headed. Then flip again, get heads again, and update once
more. How many consecutive heads before you are more than $90\%$ sure it is two-headed?
\end{exercise}

\begin{exercise}[Beta--Binomial by hand]
You hold a $\Bet(2,3)$ prior on a website's click-through rate $\theta$. You then serve the
ad to $50$ users and observe $20$ clicks. (a) Write the posterior distribution. (b) Compute
the posterior mean and the MAP estimate. (c) Explain in one sentence why they differ.
\end{exercise}

\begin{exercise}[The prior as pseudocounts]
Show that the posterior mean $\dfrac{\alpha+k}{\alpha+\beta+n}$ of a Beta--Binomial update
equals the sample proportion $k/n$ when $\alpha=\beta=0$ (the improper ``Haldane'' prior),
and equals $1/2$ regardless of data when $n=0$ and $\alpha=\beta$. Interpret both limits in
words.
\end{exercise}

\begin{exercise}[Gamma--Poisson forecast]
A call center receives calls as $\Pois(\lambda)$ per hour. With prior $\lambda\sim\Gam(3,1)$
you observe $4, 7, 5$ calls in three successive hours. (a) Give the posterior. (b) Give the
posterior mean and standard deviation. (c) The posterior-predictive number of calls in the
next hour is Negative Binomial; without deriving it, give its mean using
\eqref{ch31:postpred} and $\E[x_{\text{new}}\given D]=\E[\lambda\given D]$.
\end{exercise}

\begin{exercise}[Precisions add]
You measure a length with two instruments. The first reads $10.2$ with standard deviation
$0.5$; the second reads $10.0$ with standard deviation $0.2$. Treating the first reading as
a prior and the second as data (both normal, known variances), use
\eqref{ch31:normal-update} to find the combined posterior mean and standard deviation.
Which reading does the answer favor, and why?
\end{exercise}

\begin{exercise}[Conceptual: two interpretations]
A newspaper reports a ``$95\%$ confidence interval of $[0.41, 0.47]$ for the candidate's
vote share.'' (a) State precisely what this means to a frequentist. (b) State what a
Bayesian $95\%$ credible interval of $[0.41,0.47]$ would mean. (c) Explain which one
matches the way most readers will interpret the sentence, and what extra ingredient the
Bayesian statement required.
\end{exercise}

\begin{exercise}[Regularization as a prior]
Consider linear regression with weights $w$ and Gaussian observation noise. Show that
maximizing the posterior under the prior $w\sim\Normal(0,\sigma^2 I)$ is equivalent to
minimizing (squared error) $+\;\lambda\norm{w}^2$, and find $\lambda$ in terms of $\sigma$
and the noise variance. (Hint: take $-\log$ of likelihood times prior.)
\end{exercise}

\begin{exercise}[Proof: order does not matter]
Prove that updating a prior with data $D_1$ and then updating the result with independent
data $D_2$ yields the same posterior as updating once with the combined data $(D_1,D_2)$.
Identify exactly where you use the conditional independence of $D_1$ and $D_2$ given
$\theta$, and explain why this justifies calling yesterday's posterior today's prior.
\end{exercise}

\begin{exercise}[Sensitivity to the prior]
For the Beta--Binomial model with data $k=3$ successes in $n=4$ trials, compute the
posterior mean under three priors: $\Bet(1,1)$, $\Bet(10,10)$, and $\Bet(1,10)$. Comment on
how much each prior moves the estimate away from the raw proportion $3/4$, and relate the
size of the movement to the ``pseudocount'' strength $\alpha+\beta$ of each prior.
\end{exercise}

\begin{exercise}[MAP can mislead]
Sketch (or describe) a posterior density on $[0,1]$ that has two well-separated peaks of
equal height, say near $0.2$ and $0.8$, with a deep valley between. Explain why reporting
the MAP estimate alone would be seriously misleading here, what the posterior mean would
report, and why neither single number is an adequate summary. What would you report
instead?
\end{exercise}
